\chapter{Первый рабочий срез: Stage 1}
\label{ch:stage1}

\section{Зачем отдельная глава про Stage 1}

У больших архитектурных текстов есть повторяющаяся болезнь: чем лучше
описана система в теории, тем легче откладывать момент, когда она должна
столкнуться с реальной инженерной средой. Возникает соблазн дописать ещё
одну красивую главу, ещё одно уточнение к консолидации или ещё один будущий
модуль, прежде чем строить хоть что-то наблюдаемое.

Stage 1 нужен как лекарство именно от этого. Это не ``урезанная версия
мечты'', а минимальный рабочий срез, на котором уже видна физика системы:

\begin{itemize}\itemsep2pt
\item как событие входит в память;
\item как оно превращается в typed node;
\item как решается вопрос commit/skip;
\item как устроено извлечение;
\item как собирается контекст для LLM;
\item как вся эта внутренняя механика наблюдается и проверяется извне.
\end{itemize}

\section{Чего Stage 1 должен добиться}

На первой стадии не требуется полноценная ``живая память Рин'' во всём
богатстве книги. Требуется более скромная, но более важная вещь: сделать
внутреннюю механику памяти наблюдаемой и проверяемой.

\begin{principle}[Первый успех формулируется инженерно]
Критерий успеха Stage 1 должен звучать не как ``Рин выглядит умной'', а как
``можно руками проследить путь от события до узла памяти, от узла до retrieval
и от retrieval до контекста модели''. Пока эта прозрачность не достигнута,
разговор об ``умном поведении'' почти всегда будет смешивать качество памяти с
шумом LLM.
\end{principle}

\section{Сквозной pipeline v0}

Минимальный рабочий срез удобно описывать как цепочку:

\begin{lstlisting}[style=pseudo]
event_log             # append observed/committed/deleted
  -> candidate extraction # на Stage 1 может быть ручной
  -> typed nodes      # node_kind обязателен
  -> write/skip gate  # epistemic + salience + anti-fusion
  -> forensic index   # метаданные / provenance / время
  -> anchor retrieval
  -> local graph expansion
  -> context-packer v0 # observed / hypothesis / low_conf / ...
  -> inspector CLI
  -> synthetic tests
\end{lstlisting}

В этой цепочке важно не количество компонентов, а то, что каждый из них
имеет наблюдаемый вход и наблюдаемый выход. Если какой-то шаг ``случается
внутри магии'', Stage 1 свою задачу не выполняет.

\section{Что должно быть реализовано в первом срезе}

\textbf{1. Event log}

Нужен неизменяемый журнал событий: что пришло на вход, когда это произошло,
откуда получено и во что в итоге было превращено. Даже если candidate позже
будет отклонён, след о самом факте рассмотрения должен сохраняться. В
терминах контракта Stage 1 это append-only лог с событиями вроде
\texttt{observed}, \texttt{committed} и \texttt{deleted}; именно он делает
возможным replay и event sourcing.

\textbf{2. Typed node store}

Нужен минимальный, но уже честный store узлов и рёбер. Узел Stage 1 не обязан
содержать все будущие поля книги, но обязан нести базовую схему:

\begin{itemize}\itemsep2pt
\item \texttt{id};
\item \texttt{content};
\item \texttt{embedding};
\item \texttt{node\_kind};
\item \texttt{confidence};
\item \texttt{provenance} (или support-skeleton);
\item \texttt{time};
\item \texttt{layer}.
\end{itemize}

Остальные поля вроде \texttt{valence}, \texttt{importance}, \texttt{domains}
или \texttt{need\_to\_*} могут достраиваться позже, но \texttt{node\_kind}
обязателен с самого начала: именно он задаёт смысл \texttt{confidence}.

\begin{intuition}[Чего нельзя откладывать]
Поле \texttt{node\_kind} нельзя считать ``будущим украшением''. Без него
невозможно даже корректно интерпретировать \texttt{confidence}: степень
доверия к факту, устойчивость предпочтения и надёжность реактивного паттерна
--- разные вещи.
\end{intuition}

\textbf{3. Write/skip gate}

Нужен gate, который умеет отвечать на вопрос: следует ли вообще писать это в
память? Даже в простейшем виде он должен различать:

\begin{itemize}\itemsep2pt
\item точный дубль;
\item незначимый candidate;
\item чрезмерно разросшийся кандидат;
\item допустимый commit.
\end{itemize}

Особенно важно, чтобы Stage 1 логировал не только коммиты, но и причины
пропуска. В инженерной разбивке такой gate уже на первом срезе полезно
мыслить как тройку:

\begin{itemize}\itemsep2pt
\item \textbf{epistemic-gate:} стоит ли этому верить и с какой уверенностью;
\item \textbf{salience-gate:} значимо ли это для дальнейшей жизни памяти;
\item \textbf{anti-fusion-gate:} не пытается ли один candidate притащить
слишком много структуры разом.
\end{itemize}

Без видимых причин \texttt{skip} память слишком быстро превращается в чёрный
ящик.

\textbf{4. Forensic index}

Нужен forensic index: точный индекс по метаданным, provenance, времени и
служебным полям. Он нужен не для живости retrieval, а для диагностики,
точного поиска, удаления и реплея. Если у памяти нет такого слоя, она быстро
становится удобной для ассоциаций, но неудобной для инженерного контроля. На
первом срезе сюда же относится и работающий delete-cascade хотя бы в
минимально проверяемой форме.

\textbf{5. Anchor retrieval}

Нужен первый работающий retrieval, пусть даже без всего богатства книги. В
минимуме это означает:

\begin{itemize}\itemsep2pt
\item векторный поиск якорей;
\item ограниченное расширение окрестности на графе;
\item внятный и воспроизводимый результат.
\end{itemize}

Именно на Stage 1 впервые становится видно, что память --- это не
``дополнительный prompt'', а отдельная поисковая и структурная машина.

\textbf{6. Context-packer}

Нужен packer, который собирает память в форму, пригодную для подачи в модель.
Это ещё не финальный формат выражения, но уже важный тест: система должна
уметь не просто хранить узлы, а превращать их в рабочий контекст с
эпистемическими пометками. На первом срезе такие пометки должны быть хотя бы
базовыми: \texttt{observed}, \texttt{hypothesis}, \texttt{low\_conf} и
родственными им.

\textbf{7. Inspector}

Нужен inspector памяти: инструмент, позволяющий видеть узлы, рёбра, поля,
provenance, gate decisions и результаты retrieval. Без него Stage 1 будет
трудно не только улучшать, но даже честно диагностировать.

\textbf{8. Synthetic tests}

Нужен стенд на маленьких контролируемых графах. На этой стадии особенно
важно тестировать не ``правдоподобие диалога'', а локальные свойства
архитектуры:

\begin{itemize}\itemsep2pt
\item правильно ли работает запись;
\item не теряются ли provenance и типы;
\item даёт ли retrieval осмысленную окрестность;
\item воспроизводим ли replay из event log;
\item видны ли базовые failure modes.
\end{itemize}

\section{Что сознательно откладывается}

Stage 1 силён именно тем, что он ограничен. В эту стадию сознательно не стоит
пытаться втиснуть всё, что описано в книге. Иначе она перестанет быть первым
срезом и превратится в бесконечную ``почти готовую систему''.

На раннем этапе допустимо отложить:

\begin{itemize}\itemsep2pt
\item полную динамику настроения;
\item сложную personality-driven маршрутизацию;
\item тяжёлый fusion в полном виде, оставляя на ранней стадии только
anti-fusion-скелет и провизорные кандидаты;
\item полную TMS-ретракцию;
\item продвинутый сон и скраб;
\item гиперграфы, клеточные комплексы и другие дальние ветки;
\item тонкую governance-политику удаления и криптостирания;
\item мета-обучение операторов.
\end{itemize}

\begin{intuition}[Смысл ограничений]
Отложенные механизмы здесь не ``отменяются'', а переносятся до тех пор, пока
нижний срез не станет доверенным. Иначе система получает богатую теорию
поверх недоверенного основания.
\end{intuition}

\section{Три инженерных упрощения первой версии}

Даже внутри Stage 1 полезно заранее зафиксировать несколько упрощений, чтобы
не симулировать преждевременную полноту.

\textbf{Tombstone вместо полного забывания}

На первом проходе достаточно tombstone-подхода: узел помечается как удалённый
и исчезает из рабочей проекции, но система ещё не обязана реализовывать
полную каскадную очистку всех производных следов.

\textbf{Компоненты связности вместо полной консолидации сообществ}

На первом проходе достаточно более грубой схемы предложений к консолидации:
компоненты связности и жёсткие пороги пересечения вместо полного community
detection и дорогого graph coarsening. На раннем этапе важнее добиться
корректности typed storage и retrieval, чем сделать избыточно умное слияние. В
более инженерной формулировке: вместо тяжёлого со-консолидационного графа и
поиска модулярности достаточно простых connected components по графу
предложений с отсечением по жёсткому порогу пересечения, например в духе
LSH/MinHash.

\textbf{Простой семантический вентиль}

Stage 1 может жить с довольно грубым семантическим гейтом, если это честно
осознано и изолировано. Важна не максимальная красота фильтра, а
воспроизводимость его ошибок и возможность их позже улучшать. На уровне
roadmap это формулируется ещё конкретнее: первым проходом допустим просто
косинус кэш-эмбеддингов. Более сильная защита от омонимии через отдельный
\texttt{sense/type-gate} требует sense-слоя, которого на Stage 1 ещё нет,
поэтому здесь сознательно не обещается полная защита от омонимии.

\section{Критерии готовности}

Stage 1 можно считать состоявшимся, если выполнены следующие условия.

\begin{enumerate}\itemsep2pt
\item Событие действительно проходит путь ``лог $\to$ candidate $\to$ gate
$\to$ commit/skip'', и этот путь можно наблюдать.
\item Типы узлов и provenance сохраняются и доступны inspector'у.
\item tombstone-удаление сносит узел из рабочей проекции, и это проверяемо
inspector'ом.
\item Retrieval отрабатывает на синтетическом графе и возвращает
содержательно осмысленную окрестность.
\item Replay event log воспроизводит ту же графовую проекцию.
\item Хотя бы один базовый failure mode из батареи отказов --- например,
концентрация $\pi$ на хабах или точное удаление --- обнаруживается и
воспроизводимо чинится на стенде.
\end{enumerate}

\begin{intuition}[Как должен выглядеть типичный успех Stage 1]
На вход подаётся контролируемое событие. Затем в inspector'е видно: сам
event, candidate, причину commit или skip, созданный узел, связанные с ним
рёбра, найденные retrieval-якоря, собранный packer'ом контекст и, если нужно,
support-skeleton. После этого тот же сценарий прогоняется на synthetic graph и
даёт воспроизводимый результат. Именно такая наблюдаемость и есть настоящее
достижение первой стадии.
\end{intuition}

\begin{principle}[Первый успех --- не ``Рин помнит'', а ``физика графа наблюдаема'' ]
Первым продуктом этой стадии должна быть не ``память v0'' в рекламном
смысле, а memory inspector v0. Если первым мерилом сделать ``Рин умно
пользуется памятью в диалоге'', система почти сразу утонет в шуме языковой
модели. Поэтому успех Stage 1 состоит в том, что можно руками создать
событие, увидеть candidate, увидеть commit или skip, увидеть узел, увидеть
retrieval, увидеть provenance или support-skeleton и затем воспроизвести тот
же путь на synthetic graph.
\end{principle}

\section{Роль Stage 1 в архитектуре всей книги}

Stage 1 не является ``маленькой технической вставкой'' в конце курса. Он
играет методологическую роль для всей архитектуры памяти.

\begin{itemize}\itemsep2pt
\item Он заставляет различать красивые понятия и реально наблюдаемые
механизмы.
\item Он принуждает к дисциплине типов, provenance и gate decisions.
\item Он отделяет дефекты памяти от дефектов языковой модели.
\item Он создаёт первую твёрдую площадку, на которой дальнейшие главы можно
не только обсуждать, но и проверять.
\end{itemize}

Поэтому с точки зрения разработки Stage 1 --- не скромное начало, а первый
момент, когда архитектура памяти вообще начинает иметь право на дальнейшее
усложнение.

\begin{principle}[Дальше --- строить и ломать, а не расширять]
После фиксации этого среза следующий лучший ход --- не писать ещё одну
красивую главу, а собирать систему, прогонять synthetic tests и позволять
стенду править гипотезы. Каждый следующий красивый модуль имеет право
войти в курс только тогда, когда нижний срез уже выдерживает нагрузку.
\end{principle}
